% GENERATED FILE. Edit canonical lessons, then run npm run generate:books.
% canonical-source-hash: fnv1a64:e750b3243ecc536e
% canonical-lessons: AR-C132-sahil, AR-C132-shallal, AR-C132-kahf, AR-C132-wadi, AR-C132-tall

\chapter{Coast, Waterfall, Cave, Valley, Hill}
\label{ch:ar-coast-waterfall-cave-valley-hill}
\hlchaptermodality{132}

\begin{chapteropening}
\textbf{By the end of this chapter:} \emph{I can say a coast, a waterfall, a cave, a valley and a hill.}

Complete the last lesson of chapter 132: I can say a coast, a waterfall, a cave, a valley and a hill.
\end{chapteropening}

\section[\texorpdfstring{\ar{ساحل} (sāḥil)}{ساحل (sāḥil)}]{\ar{ساحل} (sāḥil) — a coast}

\label{lesson:AR-C132-sahil}



\begin{quote}
Before the new one: say the Arabic for a branch, then the Arabic for an island.
\end{quote}

\subsection*{You'll want to know: \ar{ساحل}}
\textbf{\ar{ساحل}} — \emph{sāḥil} — \textquotedblleft{}a coast\textquotedblright{}.

\begin{cousinweb}
\textbf{\ar{ساحل}} (\emph{sāḥil}) is built on the native Arabic root \emph{s-ḥ-l}, \textquotedblleft{}to scrape, to pare\textquotedblright{} (Lane). Lane gives \emph{sāḥil} as the shore; Swahili comes from its plural \emph{sawāḥil}, \textquotedblleft{}coasts\textquotedblright{}.
\end{cousinweb}

\begin{grammarlens}[title={What you've built}]
One more word for this chapter.
\end{grammarlens}

\subsection*{Your turn}
\begin{itemize}
\raggedright
  \item \emph{Say it:} \emph{sāḥil}
  \item \emph{Say it:} \emph{sāḥil}, once more
  \item \emph{Say it:} say \emph{sāḥil}
  \item \emph{Recall:} say the Arabic for a branch, then the Arabic for an island, then say \emph{sāḥil} again
\end{itemize}

\begin{tcolorbox}[breakable,colback=teal!4,colframe=teal!35!black,title={Before you move on}]
What does \ar{ساحل} mean? (\textquotedblleft{}A coast\textquotedblright{}.) Say it once more.
\end{tcolorbox}
\section[\texorpdfstring{\ar{شلال} (shallāl)}{شلال (shallāl)}]{\ar{شلال} (shallāl) — a waterfall}

\label{lesson:AR-C132-shallal}



\begin{quote}
Before the new one: say the Arabic for an island, then the Arabic for a coast.

\emph{Read it:} \textbf{\ar{أخذ}} in its three pieces (\emph{ʾa-kha-dha}, \textquotedblleft{}he took.\textquotedblright{})
\end{quote}

\subsection*{You'll want to know: \ar{شلال}}
\textbf{\ar{شلال}} — \emph{shallāl} — \textquotedblleft{}a waterfall\textquotedblright{}.

\begin{cousinweb}
\textbf{\ar{شلال}} (\emph{shallāl}) is built on the native Arabic root \emph{sh-l-l}, \textquotedblleft{}to drive away\textquotedblright{} (Lane). \emph{Shallāl} is a later word on this root, \textquotedblleft{}a cascade\textquotedblright{}.
\end{cousinweb}

\begin{grammarlens}[title={What you've built}]
One more word for this chapter.
\end{grammarlens}

\subsection*{Your turn}
\begin{itemize}
\raggedright
  \item \emph{Say it:} \emph{shallāl}
  \item \emph{Say it:} \emph{shallāl}, once more
  \item \emph{Say it:} say \emph{shallāl}
  \item \emph{Recall:} say the Arabic for an island, then the Arabic for a coast, then say \emph{shallāl} again
\end{itemize}

\begin{tcolorbox}[breakable,colback=teal!4,colframe=teal!35!black,title={Before you move on}]
What does \ar{شلال} mean? (\textquotedblleft{}A waterfall\textquotedblright{}.) Say it once more.
\end{tcolorbox}
\section[\texorpdfstring{\ar{كهف} (kahf)}{كهف (kahf)}]{\ar{كهف} (kahf) — a cave}

\label{lesson:AR-C132-kahf}



\begin{quote}
Before the new one: say the Arabic for a coast, then the Arabic for a waterfall.
\end{quote}

\subsection*{You'll want to know: \ar{كهف}}
\textbf{\ar{كهف}} — \emph{kahf} — \textquotedblleft{}a cave\textquotedblright{}.

\begin{cousinweb}
\textbf{\ar{كهف}} (\emph{kahf}) is built on the native Arabic root \emph{k-h-f}, \textquotedblleft{}a cave\textquotedblright{} (Lane). An old noun; Lane uses it for a cavern in a mountain.
\end{cousinweb}

\begin{grammarlens}[title={What you've built}]
One more word for this chapter.
\end{grammarlens}

\subsection*{Your turn}
\begin{itemize}
\raggedright
  \item \emph{Say it:} \emph{kahf}
  \item \emph{Say it:} \emph{kahf}, once more
  \item \emph{Say it:} say \emph{kahf}
  \item \emph{Recall:} say the Arabic for a coast, then the Arabic for a waterfall, then say \emph{kahf} again
\end{itemize}

\begin{tcolorbox}[breakable,colback=teal!4,colframe=teal!35!black,title={Before you move on}]
What does \ar{كهف} mean? (\textquotedblleft{}A cave\textquotedblright{}.) Say it once more.
\end{tcolorbox}
\section[\texorpdfstring{\ar{وادي} (wādī)}{وادي (wādī)}]{\ar{وادي} (wādī) — a valley}

\label{lesson:AR-C132-wadi}



\begin{quote}
Before the new one: say the Arabic for a waterfall, then the Arabic for a cave.
\end{quote}

\subsection*{You'll want to know: \ar{وادي}}
\textbf{\ar{وادي}} — \emph{wādī} — \textquotedblleft{}a valley\textquotedblright{}.

\begin{cousinweb}
\textbf{\ar{وادي}} (\emph{wādī}) is built on the native Arabic root \emph{w-d-y}, \textquotedblleft{}a valley, a river bed\textquotedblright{} (Lane). English \emph{wadi} comes from it.
\end{cousinweb}

\begin{grammarlens}[title={What you've built}]
One more word for this chapter.
\end{grammarlens}

\subsection*{Your turn}
\begin{itemize}
\raggedright
  \item \emph{Say it:} \emph{wādī}
  \item \emph{Say it:} \emph{wādī}, once more
  \item \emph{Say it:} say \emph{wādī}
  \item \emph{Recall:} say the Arabic for a waterfall, then the Arabic for a cave, then say \emph{wādī} again
\end{itemize}

\begin{tcolorbox}[breakable,colback=teal!4,colframe=teal!35!black,title={Before you move on}]
What does \ar{وادي} mean? (\textquotedblleft{}A valley\textquotedblright{}.) Say it once more.
\end{tcolorbox}
\section[\texorpdfstring{\ar{تل} (tall)}{تل (tall)}]{\ar{تل} (tall) — a hill}

\label{lesson:AR-C132-tall}



\begin{quote}
Before the new one: say the Arabic for a cave, then the Arabic for a valley.
\end{quote}

\subsection*{You'll want to know: \ar{تل}}
\textbf{\ar{تل}} — \emph{tall} — \textquotedblleft{}a hill\textquotedblright{}.

\begin{cousinweb}
\textbf{\ar{تل}} (\emph{tall}) is built on the native Arabic root \emph{t-l-l}, \textquotedblleft{}a hill, a mound\textquotedblright{} (Lane). English \emph{tell}, an archaeologist's mound, comes from it; Hebrew \emph{tel} is a cousin.
\end{cousinweb}

\begin{grammarlens}[title={What you've built}]
One more word for this chapter.
\end{grammarlens}

\subsection*{Your turn}
\begin{itemize}
\raggedright
  \item \emph{Say it:} \emph{tall}
  \item \emph{Say it:} \emph{tall}, once more
  \item \emph{Say it:} say \emph{tall}
  \item \emph{Recall:} say the Arabic for a cave, then the Arabic for a valley, then say \emph{tall} again
\end{itemize}

\begin{tcolorbox}[breakable,colback=teal!4,colframe=teal!35!black,title={Before you move on}]
What does \ar{تل} mean? (\textquotedblleft{}A hill\textquotedblright{}.) Say it once more.
\end{tcolorbox}
